%\textcolor{red}{NOTE: I think most of the results and code for this section are from the ACER server so they are stored in the /hdd_1/politow2/ACER_bkp/experiments/MEO.zip and smv_ood.zip files. I think\dots}

%(mv cnn definitions etc.)
%(establish mathematical definitions and explain toolkit/procedure used in detail)

\section{Multi-View OOD Detection Frameworks}
\label{sec:m1-mvcnns}

%The incorporation of multi-view data in OOD detection models is established by using 1) an off-the-shelf image segmentation tool and 2) a multi-view ML model pipeline. First, an image $x$ is segmented into two ``views'': its foreground ($x_{fg}$) and background ($x_{bg}$) view. A specific image segmentation tool (IST) for performing the segmentation is detailed in Section \ref{sec:propwork-mod-exp}. The segmentation operation results in a set of foreground object pixels in $x$ which is used as the foreground view $x_{fg}$. The background view $x_{bg}$ is the negative mask of the foreground object pixels. The original image $x$, which henceforth is equivalently denoted as the original view $x_{org}$, is included with the foreground view $x_{fg}$ and the background view $x_{bg}$ to create the resultant multi-view set for given sample image $x$ after the segmentation step, \{$x_{org}$, $x_{fg}$, $x_{bg}$\}. For convenience, the set of views is denoted as $V = \{org, fg, bg\}$ and is comprised of each individual view $v \in V$. In practice, the multiple views for any given dataset can be generated offline prior to use in training or OOD detection testing to improve efficiency. 
%A multi-view dataset $\mathcal{D}_{MV} = \{(x_{MV,i}, y_i)\}$ contains set of labels $y_i$ and a set of associated multi-view samples $x_{MV,i}$. Each multi-view sample is composed of a set of \emph{views} $x_{MV,i} = \{x_{i,0}, x_{i,1}, \dots, x_{i,V}\}$ where $V$ is the number of views. Views are images, with each view being a copy of the original image under covariate or semantic shift (including no shift, i.e. the original image itself). That is, given some augmentation function introducing covariate or semantic shift $\mathcal{A}(\cdot)$, a view is defines as $x_{i,v} = \mathcal{A}(x_i)$ for original image $x_i$ and with $v \in \mathbb{N}^V$.
The challenge in image-based \gls{ood} detection has predominantly shifted to solving the problem of detecting near-\gls{ood} data \cite{yang2022openood,winkens2020contrastive}. As previously mentioned, near-\gls{ood} is considered to be \gls{ood} data that has significant covariate and/or semantic overlap with the \gls{id} data distribution. Detection of near-\gls{ood} data relies on model learning of discriminative features representative of \gls{id} data (or classes) and subsequent rejection of spurious features \cite{ming2022impact,neuhaus2023spurious}. Existing multi-view and \gls{ood} detection works primarily focus on \emph{augmenting} data through various transformations that serve to increase the diversity of data in the \gls{id} distribution and improve model learning \cite{shorten2019survey,hendrycks2022pixmix,geng2023foreground}. However, studies investigating the behaviors of ML models and the impact of spurious features on learning have revealed patterns identifying foregrounds and backgrounds of images as critical components separating discriminative and spurious features \cite{geng2023foreground,miyai2024locoop}. There has been very little work in the multi-view literature that considers this approach to multi-view learning, i.e. using foregrounds and backgrounds as views, emphasizing the need for investigation in this direction of research. It is hypothesized that using foregrounds and backgrounds of an image as views can improve model learning of discriminative features in images, thus improving \gls{ood} detection and particularly advancing the field in performance on challenging near-\gls{ood} detection settings. This also pushes ML closer to realistic deployment in open-world settings, of interest in many applications (see Chapter \ref{diss_intro}). The hypothesis is supported by observations in medical research, positing that segmentation is a critical function in rapid object identification in humans \cite{weiner2016anatomical,vuilleumier2002multiple}.

A foreground-background multi-view \gls{ood} detection method involves a pipeline containing two components: 1) a foreground-background (FgBg) \gls{ist} $f_{IST}$, and 2) a multi-view model $f_{MV}$. FgBg segmentation is an operation $\mathcal{K}$ that takes as input an image $x$ and outputs a set of foreground pixels which is used as the foreground view $x_{fg}$. That is, given that $x$ is an image, $x_{fg} = \{(\text{pixel},i,j) | (i,j) \in \mathcal{K}(x)\}$. The background view $x_{bg}$ is the result of using the foreground pixels as a negative mask on the original image $x$, $x_{bg} = x \setminus x_{fg}$. Note that here the set definitions with $\mathcal{K}$ are the result of the aforementioned view-generating augmentation function $\mathcal{B}$ mentioned in previous multi-view work (see Chapter \ref{diss_bkgd} Section \ref{sec:bkgd-mvp}). Multi-view samples are then constructed using the original image, henceforth referred to as the original view $x_{org}$, and the foreground and background views $x_{MV} = \{x_{org}, x_{fg}, x_{bg}\}$. A shorthand is introduced here denoting the set of views as $V = \{org, fg, bg\}$, comprised of each individual view $v\in V$. In practice, the FgBg views for any dataset are generated offline with the \gls{ist} prior to use in training or testing for convenience. %Analysis is provided later (see Section \ref{}) for evaluation of method complexities and relevant resource consumption of all components involved in proposed methods.

%Second, each of the three views is processed using a multi-view image processing ML model. For the subsequent proposed works, the model is a multi-view convolutional NN (CNN), or MVCNN \cite{su2015multi}. An MVCNN is composed of a set of CNN feature extractors $\Phi_{V} = \{\phi_{org}, \phi_{fg}, \phi_{bg}\}$, one for each input view $v \in V$, and a fusion module $\mathcal{G}$. The CNN feature extractors are responsible for coalescing input image data into a feature vector representative of class-relevant image information, while the fusion module is responsible for regulating how much information should be included in the final model feature vector from each of its input view feature vectors. A CNN feature extractor $\phi_{v}$ takes as input its corresponding view $x_{v}$ and outputs a $p$-dimensional feature vector $z_{v} \in \mathbb{R}^{p}$, $\phi_{v}: \mathcal{X} \mapsto \mathbb{R}^{p}$. The CNN feature extractors $\Phi_{V}$ output a set of \textbf{view-specific} feature vectors \{$z_{org}$, $z_{fg}$, $z_{bg}$\} which are then passed to a fusion module $\mathcal{G}$. A fusion module $\mathcal{G}$ takes as input feature vectors from each view $v$ and outputs a combination of feature $Z_{mv}$ and prediction logits $\hat{Y}_{mv}$. %\textcolor{blue}{do we need more detail on fusion module and in general here?} 
%Details on the MVCNNs and fusion modules used in the presented works for OOD detection are provided in Sections \ref{sec:propwork-mod-propmeth} and \ref{sec:propwork-meo-propmeth}. An example of such a pipeline is given in Figure \ref{fig:propwork-mod-diagram}. %\textcolor{blue}{maybe add some sentences here walking through the diagram and matching key features to described components}
%Multi-view models process multi-view datasets by taking each view in a multi-view sample as input and calculating some representative representation coalescing information across the views as output. The multi-view model $f_{MV}$ may be considered as a set of individual-view models $f_{MV} = \{f_0, f_1, \dots, f_V\}$. Each view in a multi-view sample $x_{i,v}$ is processed by its corresponding individual-view model $f_v$. Note that each individual-view model can be the same, and thus shared. Outputs from each individual-view model are then coalesced in an operation referred to as ``fusion'', resulting in a collective output from the multi-view model $f_{MV}(x_{MV,i}) = \mathcal{U}(\{f_0(x_{i,0}),f_1(x_{i,1}),\dots,f_V(x_{i,V})\})$ for fusion operation $\mathcal{U}$.
FgBg multi-view models are composed of a set of image processing models and a fusion component handling the fusion operation (these terms will be used interchangeably), $\mathcal{U}$. The image processing models can be decomposed into several sets, each of which is equivalent to a full image processing model. Specifically, individual view model $f_v$ is decomposed into a feature extractor $\phi_v$ and a classification layer weight matrix $W_v$ with optional bias $b_v$. Thus, the multi-view model becomes a decomposed multi-view feature extractor set $\phi_{MV} = \{\phi_v | v \in V\}$, multi-view classification weight matrix set $W_{MV} = \{W_v | v \in V\}$, and optional multi-view classification bias set $b_{MV} = \{b_v | v \in V\}$. Individual view features $z_v = \phi_v(x)$ and predicted logits $\hat{y}_v = W_v z_v + b_v$ are represented by the multi-view sets $Z_{MV} = \{z_v | v \in V\}$ and $\hat{Y}_{MV} = \{\hat{y}_v | v \in V\}$. The fusion component $\mathcal{U}$ outputs a fused feature $z_{\mathcal{U}}$ and fused predicted logits $\hat{y}_{\mathcal{U}} = W_{\mathcal{U}} z_{\mathcal{U}} + b_{\mathcal{U}}$ for fused classification weight matrix $W_{\mathcal{U}}$ and optional bias $b_{\mathcal{U}}$.

%Of note is that using multiple views of data \emph{is not equivalent} to using additional training data. Instead, multi-view data is a sophisticated utilization of given data specifically targeted towards learning and exploiting salient discriminative features in image data. In this case, segmentation into foreground and background views can be seen as parallel to data augmentation, exposing models to additional information using a more diverse set of images. However, the proposed method does not add additional information to any given data sample, instead only enabling better separation of information. In all presented evaluations, the same training and test datasets are used to train and evaluate the proposed methods as are used by the compared baselines.}
Of note is that using multi-view data is \textit{not equivalent} to simply using additional training data. The views $x_{MV}$ resulting from augmentations on any given input $x$ provide covariate and/or semantic information highly correlated \textit{with respect to} the original sample. As such, processing views in a multi-view framework exploits this contextual correlation while additional data foregoes this in favor of naive learning. This approach, that is of using additional data, has been shown theoretically to be insufficient in solving the \gls{ood} detection problem \cite{fang2022out}. It is further demonstrated in this dissertation that multi-view processing is necessary to exploit $x_{MV}$ for improved OOD detection (see Sections \ref{sec:propwork-mod-exp}, \ref{sec:propwork-meo-exp}, \ref{sec:propwork-scod-exp}, and Chapter \ref{diss_casod} Section \ref{sec:propwork-casod-exp}).
